%% =============================================================================
%% Markov Chain Monte Carlo Methods in MechanicsDSL — Technical Whitepaper
%% =============================================================================
%% Compiled with: pdflatex mcmc_whitepaper.tex  (×2 for references)
%% =============================================================================

\documentclass[11pt,letterpaper,twocolumn]{article}

%% ── Geometry & Layout ────────────────────────────────────────────────────────
\usepackage[
  top=0.85in, bottom=1in, left=0.65in, right=0.65in,
  columnsep=0.30in
]{geometry}

%% ── Typography ───────────────────────────────────────────────────────────────
\usepackage[T1]{fontenc}
\usepackage[sc,osf]{mathpazo}            % Palatino body + math (unified serif)
\usepackage[scaled=0.88]{inconsolata}     % Matching monospace
\usepackage{microtype}
\usepackage{setspace}
\setstretch{1.08}
\usepackage[shortlabels]{enumitem}
\setlist{nosep,leftmargin=*}              % Tighter, consistent lists

%% ── Mathematics ──────────────────────────────────────────────────────────────
\usepackage{amsmath,amssymb,amsthm,mathtools}
\usepackage{bm}

%% ── Figures & Tables ─────────────────────────────────────────────────────────
\usepackage{graphicx}
\usepackage{booktabs}
\usepackage{array}
\usepackage{multirow}
\usepackage{float}

%% ── Code Listings ────────────────────────────────────────────────────────────
\usepackage{listings}
\usepackage{xcolor}

\definecolor{codegreen}{rgb}{0.15,0.55,0.15}
\definecolor{codegray}{rgb}{0.50,0.50,0.50}
\definecolor{codepurple}{rgb}{0.50,0.00,0.50}
\definecolor{backcolour}{rgb}{0.97,0.97,0.97}
\definecolor{framecolour}{rgb}{0.82,0.82,0.82}

\lstdefinestyle{mechanicsstyle}{
  backgroundcolor=\color{backcolour},
  commentstyle=\color{codegreen}\itshape,
  keywordstyle=\color{codepurple}\bfseries,
  stringstyle=\color{red!60!black},
  basicstyle=\ttfamily\small,
  breaklines=true,
  frame=single,
  rulecolor=\color{framecolour},
  numbers=left,
  numberstyle=\scriptsize\color{codegray},
  numbersep=5pt,
  tabsize=4,
  showstringspaces=false,
  captionpos=b,
  aboveskip=8pt,
  belowskip=6pt,
}
\lstset{style=mechanicsstyle}

%% ── Hyperlinks ───────────────────────────────────────────────────────────────
\usepackage[
  colorlinks=true,
  linkcolor=blue!70!black,
  citecolor=blue!70!black,
  urlcolor=blue!60!black,
  bookmarksnumbered=true
]{hyperref}

%% ── Algorithms ───────────────────────────────────────────────────────────────
\usepackage[ruled,vlined,linesnumbered]{algorithm2e}

%% ── Section Styling ──────────────────────────────────────────────────────────
\usepackage{titlesec}
\titleformat{\section}{\normalfont\large\bfseries\scshape}{\thesection.}{0.5em}{}
\titleformat{\subsection}{\normalfont\normalsize\bfseries}{\thesubsection}{0.5em}{}
\titleformat{\subsubsection}{\normalfont\small\bfseries\itshape}{\thesubsubsection}{0.5em}{}
\titlespacing*{\section}{0pt}{14pt plus 2pt minus 2pt}{6pt plus 1pt}
\titlespacing*{\subsection}{0pt}{10pt plus 2pt minus 1pt}{4pt plus 1pt}
\titlespacing*{\subsubsection}{0pt}{8pt plus 1pt minus 1pt}{3pt plus 1pt}

%% ── Header / Footer ─────────────────────────────────────────────────────────
\usepackage{fancyhdr}
\pagestyle{fancy}
\fancyhf{}
\fancyhead[L]{\small\textsc{MechanicsDSL}}
\fancyhead[R]{\small\textsc{MCMC Technical Whitepaper}}
\fancyfoot[C]{\small\thepage}
\renewcommand{\headrulewidth}{0.4pt}
\renewcommand{\footrulewidth}{0pt}

%% ── Custom Commands ──────────────────────────────────────────────────────────
\newcommand{\code}[1]{\texttt{#1}}
\newcommand{\pkg}[1]{\textit{#1}}

\theoremstyle{definition}
\newtheorem{definition}{Definition}
\newtheorem{proposition}{Proposition}

%% =============================================================================
%%  DOCUMENT
%% =============================================================================
\begin{document}

%% ── Title Block ──────────────────────────────────────────────────────────────
\twocolumn[{%
\begin{@twocolumnfalse}
\begin{center}

{\LARGE\bfseries\scshape
Markov Chain Monte Carlo Methods\\[3pt]
in MechanicsDSL}\\[14pt]

{\large\itshape A Technical Whitepaper}\\[10pt]

{\normalsize
MechanicsDSL Development Team\\[2pt]
\texttt{https://github.com/MechanicsDSL/mechanicsdsl}}\\[6pt]

{\small Version 1.0\enspace---\enspace March 2026}

\end{center}

\vspace{6pt}
\hrule height 0.5pt
\vspace{8pt}

\begin{abstract}
\noindent
This whitepaper provides a formal, comprehensive treatment of the Markov Chain Monte Carlo (MCMC) and Monte Carlo (MC) methods implemented within the \pkg{MechanicsDSL} computational physics framework. \pkg{MechanicsDSL} integrates MCMC across three distinct subsystems: (1)~a \textbf{Bayesian parameter inference engine} using the Metropolis--Hastings algorithm for solving inverse problems in classical mechanics, (2)~a \textbf{Metropolis Monte Carlo simulator} for statistical-mechanics lattice models, and (3)~a \textbf{GPU\nobreakdash-accelerated batch Monte Carlo} code generator targeting CUDA-capable hardware for massively parallel parameter sweeps. We detail the mathematical foundations, algorithmic design, implementation architecture, and computational considerations of each subsystem, situating them within the broader \pkg{MechanicsDSL} compiler pipeline.
\end{abstract}

\vspace{4pt}
\hrule height 0.5pt
\vspace{14pt}
\end{@twocolumnfalse}
}]

%% =====================================================================
\section{Introduction}\label{sec:intro}
%% =====================================================================

\subsection{Motivation}

Deterministic simulation of physical systems---while powerful---provides only point estimates of system behaviour for fixed parameter values.  In practice, physical parameters carry uncertainty from measurement error, model inadequacy, and inherent stochasticity.  Markov Chain Monte Carlo methods provide a principled framework for:

\begin{itemize}
  \item \textbf{Inverse problems:} Inferring unknown physical parameters from noisy observational data via Bayesian posterior sampling.
  \item \textbf{Statistical mechanics:} Sampling equilibrium configurations of many-body systems whose exact partition functions are intractable.
  \item \textbf{Uncertainty propagation:} Quantifying how parameter uncertainty maps to prediction uncertainty through Monte Carlo ensemble simulation.
\end{itemize}

\noindent
\pkg{MechanicsDSL} addresses all three use cases through tightly integrated MCMC subsystems that share the framework's symbolic--numeric compiler pipeline.

\subsection{Scope}

Table~\ref{tab:subsystems} summarises the three MCMC/MC implementations covered by this whitepaper.

\begin{table}[H]
\centering
\caption{MCMC subsystems in \pkg{MechanicsDSL}.}\label{tab:subsystems}
\small
\setlength{\tabcolsep}{4pt}
\begin{tabular}{@{}p{1.9cm}p{2.1cm}p{2.6cm}@{}}
\toprule
\textbf{Subsystem} & \textbf{Algorithm} & \textbf{Module} \\
\midrule
Bayesian inference  & Metropolis--Hastings & \code{inverse.uncertainty} \\
Lattice MC          & Single-site Metropolis & \code{domains.statistical} \\
GPU batch MC        & Parallel RK4 ensembles & \code{codegen.cuda} \\
\bottomrule
\end{tabular}
\end{table}

%% =====================================================================
\section{Mathematical Foundations}\label{sec:math}
%% =====================================================================

\subsection{The Metropolis--Hastings Algorithm}

The Metropolis--Hastings (MH) algorithm constructs a Markov chain whose stationary distribution is a target distribution~$\pi(\bm{\theta})$.  Given a current state~$\bm{\theta}^{(t)}$:

\begin{enumerate}
  \item \textbf{Propose} $\bm{\theta}' \sim q(\bm{\theta}' \mid \bm{\theta}^{(t)})$.
  \item \textbf{Compute} the acceptance ratio
  \begin{equation}\label{eq:mh-ratio}
    \alpha = \min\!\Bigg(1,\;
      \frac{\pi(\bm{\theta}')\;q(\bm{\theta}^{(t)} \mid \bm{\theta}')}
           {\pi(\bm{\theta}^{(t)})\;q(\bm{\theta}' \mid \bm{\theta}^{(t)})}
    \Bigg).
  \end{equation}
  \item \textbf{Accept} $\bm{\theta}^{(t+1)} = \bm{\theta}'$ with probability~$\alpha$; otherwise set $\bm{\theta}^{(t+1)} = \bm{\theta}^{(t)}$.
\end{enumerate}

\noindent
For symmetric proposals $q(\bm{\theta}' \mid \bm{\theta}) = q(\bm{\theta} \mid \bm{\theta}')$ (e.g.\ Gaussian random walk), the ratio simplifies to the \emph{Metropolis ratio}:
\begin{equation}\label{eq:metropolis-ratio}
  \alpha = \min\!\left(1,\;\frac{\pi(\bm{\theta}')}{\pi(\bm{\theta}^{(t)})}\right).
\end{equation}

\begin{proposition}[Detailed Balance]
The transition kernel~$T$ satisfies detailed balance with respect to~$\pi$:
\begin{equation}
  \pi(\bm{\theta})\,T(\bm{\theta} \to \bm{\theta}')
  = \pi(\bm{\theta}')\,T(\bm{\theta}' \to \bm{\theta}),
\end{equation}
guaranteeing convergence to~$\pi$ as the stationary distribution.
\end{proposition}

\subsection{Bayesian Inference Framework}

For parameter inference the target distribution is the posterior
\begin{equation}\label{eq:posterior}
  \pi(\bm{\theta} \mid \mathcal{D})
    \propto \mathcal{L}(\mathcal{D} \mid \bm{\theta})\;p(\bm{\theta}),
\end{equation}
where $\mathcal{L}(\mathcal{D} \mid \bm{\theta})$ is the likelihood and $p(\bm{\theta})$ is the prior.  The log-posterior decomposes as
\begin{equation}\label{eq:log-posterior}
  \ln\pi(\bm{\theta} \mid \mathcal{D})
    = \ln\mathcal{L}(\mathcal{D} \mid \bm{\theta})
    + \ln p(\bm{\theta})
    + \mathrm{const.}
\end{equation}

\subsection{Statistical Mechanics and the Boltzmann Distribution}

In statistical mechanics the target is the canonical distribution at inverse temperature $\beta = 1/(k_B T)$:
\begin{equation}\label{eq:boltzmann}
  P(\mathbf{s}) = \frac{1}{Z}\exp\!\big({-\beta\,H(\mathbf{s})}\big),
  \quad
  Z = \sum_{\mathbf{s}} \exp\!\big({-\beta\,H(\mathbf{s})}\big).
\end{equation}

For the Ising Hamiltonian
\begin{equation}\label{eq:ising-H}
  H = -J\!\sum_{\langle i,j\rangle} s_i s_j
      - h\!\sum_i s_i,
\end{equation}
a single-spin flip at site~$i$ yields the energy difference
\begin{equation}\label{eq:delta-E}
  \Delta E = 2J\,s_i \sum_{j\in\mathcal{N}(i)} s_j + 2h\,s_i,
\end{equation}
accepted with probability $\min(1,\,e^{-\beta\,\Delta E})$.

%% =====================================================================
\section{Subsystem~I: Bayesian Parameter Inference}\label{sec:bayesian}
%% =====================================================================

\subsection{Architecture}

The Bayesian inference engine resides in \code{inverse/uncertainty.py} and is organised around two primary classes:

\begin{lstlisting}[language=Python,caption={Public API of the uncertainty module.}]
class UncertaintyQuantifier:
    def mcmc_inference(...)   -> MCMCResult
    def bootstrap_uncertainty(...) -> BootstrapResult
\end{lstlisting}

\noindent
The \code{UncertaintyQuantifier} is initialised with a \code{PhysicsCompiler} instance, coupling the MCMC sampler directly to the symbolic--numeric simulation pipeline.  Each likelihood evaluation:

\begin{enumerate}
  \item sets parameter values on the compiled simulator,
  \item integrates the derived equations of motion,
  \item interpolates predictions to observation times via cubic splines (\code{scipy.interpolate.interp1d}),
  \item computes residuals against observed data.
\end{enumerate}

\subsection{Likelihood Function}

The implementation assumes a Gaussian observation model:
\begin{equation}\label{eq:loglik}
  \ln\mathcal{L}
    = -\frac{1}{2}\sum_{i=1}^{N}\!\left(\frac{y_i^{\mathrm{obs}} - y_i^{\mathrm{pred}}(\bm{\theta})}{\sigma}\right)^{\!2}
      - N\ln\!\big(\sigma\sqrt{2\pi}\big),
\end{equation}
where $y_i^{\mathrm{pred}}(\bm{\theta})$ are obtained by running the compiled ODE simulation.  The noise standard deviation~$\sigma$ is estimated as $0.1 \times \mathrm{std}(\mathcal{D})$ when not provided.  Failed simulations return $\ln\mathcal{L} = -\infty$, automatically rejecting unphysical parameter proposals.

\subsection{Prior Distributions}

Three prior families are supported (Table~\ref{tab:priors}).

\begin{table}[H]
\centering
\caption{Supported prior distributions.}\label{tab:priors}
\small
\setlength{\tabcolsep}{4pt}
\begin{tabular}{@{}llp{3.2cm}@{}}
\toprule
\textbf{Prior} & \textbf{Params} & \textbf{Log-density} \\
\midrule
Uniform   & $(a,b)$         & $-\ln(b{-}a)$ if $\theta\!\in\![a,b]$ \\
Normal    & $(\mu,\sigma)$  & $-\frac{1}{2}\bigl(\frac{\theta-\mu}{\sigma}\bigr)^{2}$ \\
Log-normal& $(\mu,\sigma)$  & $-\frac{1}{2}\bigl(\frac{\ln\theta-\mu}{\sigma}\bigr)^{2} - \ln\theta$ \\
\bottomrule
\end{tabular}
\end{table}

\noindent
When no prior is specified, the system defaults to a wide uniform prior $\mathcal{U}(0.01\,\theta_0,\;100\,\theta_0)$, where $\theta_0$ is the current parameter value.

\subsection{Proposal Mechanism}

The proposal distribution is an isotropic Gaussian random walk:
\begin{equation}\label{eq:proposal}
  \theta'_j = \theta^{(t)}_j + \epsilon_j,
  \qquad
  \epsilon_j \sim \mathcal{N}\!\big(0,\;s^2\,|\theta^{(t)}_j|^2\big),
\end{equation}
where $s$ is the \code{proposal\_scale} hyperparameter (default: 0.1).  The scaling is \emph{relative to the current parameter magnitude}, ensuring scale-invariant exploration across parameters of different orders of magnitude.

Since $q(\bm{\theta}' \mid \bm{\theta}) = q(\bm{\theta} \mid \bm{\theta}')$ for symmetric Gaussian proposals, the log-acceptance ratio becomes:
\begin{equation}
  \ln\alpha = \big[\ln p(\bm{\theta}') + \ln\mathcal{L}(\mathcal{D}\mid\bm{\theta}')\big]
            - \big[\ln p(\bm{\theta}^{(t)}) + \ln\mathcal{L}(\mathcal{D}\mid\bm{\theta}^{(t)})\big].
\end{equation}
Accept if $\ln U < \ln\alpha$ where $U \sim \mathrm{Uniform}(0,1)$.

\subsection{Chain Management}

\begin{itemize}
  \item \textbf{Burn-in:} The first \code{n\_burn} samples (default: 1\,000) are discarded.  The chain is initialised at the prior mean (normal/log-normal) or the midpoint (uniform).
  \item \textbf{Post-burn-in:} \code{n\_samples} samples (default: 10\,000) form the posterior approximation.
  \item \textbf{Storage:} The full chain is stored as a NumPy array of shape $(N_{\mathrm{samples}},\,d)$.
\end{itemize}

\subsection{Output Statistics}

The \code{MCMCResult} dataclass encapsulates the quantities in Table~\ref{tab:mcmc-output}.

\begin{table}[H]
\centering
\caption{Fields of \code{MCMCResult}.}\label{tab:mcmc-output}
\small
\setlength{\tabcolsep}{4pt}
\begin{tabular}{@{}lp{3.8cm}@{}}
\toprule
\textbf{Field} & \textbf{Description} \\
\midrule
\code{samples}            & Full posterior matrix $(N \!\times\! d)$ \\
\code{param\_names}       & Parameter identifiers \\
\code{acceptance\_rate}   & Fraction of accepted proposals \\
\code{log\_likelihood}    & Log-posterior trace \\
\code{mean}               & Posterior means $\hat{\bm{\theta}}$ \\
\code{std}                & Posterior standard deviations \\
\code{credible\_intervals}& 95\% highest-density intervals \\
\bottomrule
\end{tabular}
\end{table}

\subsection{The Inverse Problem Pipeline}

The MCMC engine is one component of a three-tier architecture (Figure~\ref{fig:pipeline}).

\begin{figure}[H]
\centering
\fbox{\parbox{0.92\linewidth}{\small
\textbf{Tier 1 --- Point Estimation} \\
\code{ParameterEstimator}: L-BFGS-B, Nelder--Mead, differential evolution, grid search.  Yields $\hat{\bm{\theta}} = \arg\max \mathcal{L}(\mathcal{D}\mid\bm{\theta})$.\\[4pt]
\textbf{Tier 2 --- Uncertainty Quantification} \\
\code{UncertaintyQuantifier}: MCMC posterior $\pi(\bm{\theta}\mid\mathcal{D})$, bootstrap confidence intervals.\\[4pt]
\textbf{Tier 3 --- Sensitivity Analysis} \\
\code{SensitivityAnalyzer}: Sobol variance decomposition (Saltelli sampling), Morris screening, local finite-difference sensitivity.
}}
\caption{Three-tier inverse problem pipeline.}\label{fig:pipeline}
\end{figure}

\noindent
A typical workflow: (1)~fit parameters via optimisation, (2)~quantify uncertainty via MCMC using fitted values as prior centres, (3)~assess sensitivity via Sobol analysis to determine dominant parameters.

%% =====================================================================
\section{Subsystem~II: Metropolis Monte Carlo for the Ising Model}\label{sec:ising}
%% =====================================================================

\subsection{Lattice Representation}

The spin configuration $\{s_i\}$ is stored as a NumPy integer array with periodic boundary conditions (Table~\ref{tab:lattice}).

\begin{table}[H]
\centering
\caption{Lattice storage by dimension.}\label{tab:lattice}
\small
\setlength{\tabcolsep}{4pt}
\begin{tabular}{@{}lll@{}}
\toprule
\textbf{Dim.} & \textbf{Shape} & \textbf{Sites $N$} \\
\midrule
1D & $(L,)$         & $L$     \\
2D & $(L,L)$        & $L^2$   \\
3D & $(L,L,L)$      & $L^3$   \\
\bottomrule
\end{tabular}
\end{table}

\noindent
Spins take values $s_i \in \{-1,+1\}$ stored as \code{int8}.  Periodic boundary conditions are implemented via \code{numpy.roll}, yielding a torus topology.

\subsection{Single-Site Metropolis Update}

Algorithm~\ref{alg:metropolis} formalises the \code{metropolis\_step()} method for the 2D case.

\begin{algorithm}[H]
\SetAlgoLined
\KwIn{Lattice $\mathbf{s}$ of size $L\times L$, coupling $J$, field $h$, $\beta$}
\KwOut{Updated lattice $\mathbf{s}$, acceptance flag}
Select random site $(i,j)$, \; $i,j \in \{0,\ldots,L{-}1\}$\;
$\Sigma_{\mathrm{nn}} \gets s_{i-1,j} + s_{i+1,j} + s_{i,j-1} + s_{i,j+1}$
\quad\tcp*{periodic}
$\Delta E \gets 2J\,s_{i,j}\,\Sigma_{\mathrm{nn}} + 2h\,s_{i,j}$\;
\eIf{$\Delta E \le 0$ \textbf{or} $U < \exp(-\beta\,\Delta E)$}{
  $s_{i,j} \gets -s_{i,j}$\;
  \Return{accepted}\;
}{
  \Return{rejected}\;
}
\caption{Metropolis single-site update (2D).}\label{alg:metropolis}
\end{algorithm}

\noindent
The energy-difference formula exploits locality: only nearest-neighbour interactions contribute, yielding $O(1)$ cost per trial regardless of system size.

\subsection{Monte Carlo Sweep}

A \emph{sweep} consists of $N = L^d$ single-site trials (one attempted flip per site on average), returning:

\begin{itemize}
  \item \textbf{Acceptance rate:} fraction of proposed flips accepted.
  \item \textbf{Total energy:} computed via vectorised \code{numpy.roll} in $O(N)$.
\end{itemize}

\subsection{Physical Observables}

\begin{table}[H]
\centering
\caption{Observable quantities.}\label{tab:observables}
\small
\setlength{\tabcolsep}{4pt}
\begin{tabular}{@{}p{1.6cm}p{2.4cm}p{2.6cm}@{}}
\toprule
\textbf{Observable} & \textbf{Formula} & \textbf{Method} \\
\midrule
Energy      & Eq.~\eqref{eq:ising-H}                  & \code{energy()} \\
Magn.       & $M = \sum_i s_i$                        & \code{magnetization()} \\
Magn.~dens. & $m = M/N$                               & \code{magnetization\_density()} \\
$T_c$ (2D)  & $2J / \ln(1{+}\sqrt{2})$                & \code{critical\_temperature\_2d()} \\
\bottomrule
\end{tabular}
\end{table}

\noindent
The 2D critical temperature $T_c \approx 2.269\,J/k_B$ is the exact Onsager solution~\cite{onsager1944}, used for validation.

\subsection{Computational Scaling}

Performance benchmarks confirm the expected $O(L^2)$ scaling for 2D sweeps.  The test suite validates that timing ratios between successive lattice sizes $L \in \{10, 20, 40\}$ fall within $[2, 8]$, consistent with the theoretical quadratic single-sweep complexity.

%% =====================================================================
\section{Subsystem~III: GPU-Accelerated Batch Monte Carlo}\label{sec:cuda}
%% =====================================================================

\subsection{Architecture}

The CUDA code generator (\code{codegen/cuda.py}) produces GPU kernels for \emph{embarrassingly parallel} Monte Carlo simulations, where many independent replicas are simulated simultaneously with varied parameters or initial conditions.

\subsection{Batch Simulation Design}

The \code{generate\_batch\_simulation()} method emits a complete CUDA program.  The key insight is that each GPU thread integrates a completely independent copy of the equations of motion---derived symbolically from the DSL specification---with its own state vector.  This produces \emph{perfect parallel speedup} with no synchronisation overhead.

\begin{figure}[H]
\centering
\fbox{\parbox{0.92\linewidth}{\small
\textbf{Host (CPU):}
\begin{enumerate}\setlength\itemsep{0pt}
  \item Allocate $B \times d_s$ state array
  \item Initialise with random perturbations
  \item Transfer to GPU
  \item Time-stepping: launch \code{batch\_rk4\_kernel<<<B',T>>>}
\end{enumerate}
\textbf{Device (GPU):}
\begin{itemize}\setlength\itemsep{0pt}
  \item Each thread $\to$ one independent simulation
  \item Thread-local RK4 integration
  \item Symbolic EOM from DSL compilation
  \item No inter-thread communication
\end{itemize}
}}
\caption{GPU batch Monte Carlo architecture.}\label{fig:gpu-arch}
\end{figure}

\subsection{Equations of Motion on GPU}

The CUDA device function \code{compute\_derivatives()} is generated from SymPy symbolic expressions via the pipeline:
\[
  \text{DSL} \;\to\; \mathcal{L} \;\to\; \text{E--L EOM} \;\to\; \text{SymPy} \;\xrightarrow{\code{cxxcode()}} \;\text{CUDA}
\]
All transcendental functions (\code{sin}, \code{cos}, etc.)\ map directly to CUDA \code{libm} device intrinsics.

\subsection{Integration Scheme}

Each GPU thread performs a local 4th-order Runge--Kutta step:
\begin{equation}\label{eq:rk4}
  \mathbf{y}_{n+1} = \mathbf{y}_n + \frac{\Delta t}{6}\big(\mathbf{k}_1 + 2\mathbf{k}_2 + 2\mathbf{k}_3 + \mathbf{k}_4\big),
\end{equation}
with stages
\begin{align}
  \mathbf{k}_1 &= f(t_n,\;\mathbf{y}_n), \notag\\
  \mathbf{k}_2 &= f\!\big(t_n+\tfrac{\Delta t}{2},\;\mathbf{y}_n+\tfrac{\Delta t}{2}\mathbf{k}_1\big), \notag\\
  \mathbf{k}_3 &= f\!\big(t_n+\tfrac{\Delta t}{2},\;\mathbf{y}_n+\tfrac{\Delta t}{2}\mathbf{k}_2\big), \notag\\
  \mathbf{k}_4 &= f\!\big(t_n+\Delta t,\;\mathbf{y}_n+\Delta t\,\mathbf{k}_3\big). \label{eq:rk4-stages}
\end{align}
All intermediate vectors are allocated in per-thread registers/local memory, eliminating shared-memory contention.

\subsection{Thread Mapping}

\begin{equation}
  \text{blocks} = \left\lceil \frac{B}{256} \right\rceil, \quad
  \text{threads} = 256.
\end{equation}
Thread~\code{idx} accesses states \code{[idx$\cdot d_s$ \ldots\ (idx+1)$\cdot d_s$)}.

\subsection{Monte Carlo Application Patterns}

The batch infrastructure supports:
\begin{enumerate}
  \item \textbf{Parameter sweeps:} vary physical parameters across the batch.
  \item \textbf{MC uncertainty propagation:} sample parameters from uncertainty distributions, run ensemble, compute output statistics.
  \item \textbf{Sensitivity analysis:} complement CPU-based Sobol analysis with GPU-accelerated evaluations.
\end{enumerate}

\noindent
The example in \code{gpu\_batch\_example.py} demonstrates generating code for 10\,000 parallel pendulum simulations:

\begin{lstlisting}[language=Python,caption={Batch Monte Carlo code generation.}]
gen = CudaGenerator(
    system_name="monte_carlo_pendulum",
    coordinates=['theta'],
    parameters={'m':1.0, 'l':1.0, 'g':9.81},
    batch_size=10000,
    use_cublas=True,
    compute_capability="70"
)
\end{lstlisting}

\subsection{Optional cuBLAS Integration}

For systems requiring linear algebra (constrained dynamics, mass-matrix inversions), the generator emits cuBLAS helpers (Table~\ref{tab:cublas}).

\begin{table}[H]
\centering
\caption{cuBLAS helper functions.}\label{tab:cublas}
\small
\setlength{\tabcolsep}{4pt}
\begin{tabular}{@{}lll@{}}
\toprule
\textbf{Operation} & \textbf{cuBLAS} & \textbf{Use Case} \\
\midrule
Mat--vec mult. & \code{Dgemv} & Mass matrix solve \\
Dot product    & \code{Ddot}  & Energy comp. \\
Vector norm    & \code{Dnrm2} & Convergence \\
Mat--mat mult. & \code{Dgemm} & Jacobians \\
\bottomrule
\end{tabular}
\end{table}

%% =====================================================================
\section{Cross-Cutting Analysis}\label{sec:cross}
%% =====================================================================

\subsection{Convergence Diagnostics}

The \code{MCMCResult} structure provides fundamental convergence diagnostics:

\begin{itemize}
  \item \textbf{Acceptance rate.} Roberts et al.~\cite{roberts1997} establish an optimal rate of ${\sim}\,0.234$ for random-walk Metropolis in high dimensions; the \code{proposal\_scale} parameter provides direct control.
  \item \textbf{Trace plots.} Via \code{get\_chain(param)} for marginal chain inspection.
  \item \textbf{Log-posterior trace.} Monitors mixing and stationarity.
  \item \textbf{95\% credible intervals.} Computed from the 2.5th and 97.5th percentiles of the marginal posteriors.
\end{itemize}

\subsection{Error Handling and Robustness}

\begin{itemize}
  \item \textbf{Bayesian inference:} Failed simulations (ODE divergence, interpolation failure) return $\ln\mathcal{L} = -\infty$, ensuring unphysical parameter regions are excluded from the posterior.
  \item \textbf{Ising model:} The Metropolis criterion naturally handles both energy-lowering ($\Delta E \le 0$, always accepted) and energy-raising ($\Delta E > 0$, thermally activated) transitions.
  \item \textbf{CUDA batch:} The \code{CUDA\_CHECK} macro wraps all runtime API calls, converting errors to diagnostics with file/line information.
\end{itemize}

\subsection{Bootstrap: A Frequentist Complement}

The \code{UncertaintyQuantifier} also provides frequentist uncertainty via bootstrap resampling:

\begin{itemize}
  \item \textbf{Nonparametric:} resample $(t_i, y_i)$ with replacement, refit, aggregate.
  \item \textbf{Parametric:} add synthetic noise, refit, aggregate.
\end{itemize}

\noindent
This dual Bayesian/frequentist approach enables cross-validation of uncertainty estimates.

%% =====================================================================
\section{Computational Complexity}\label{sec:complexity}
%% =====================================================================

\subsection{Bayesian MCMC Inference}

Per sample, the dominant cost is the forward simulation:
\begin{equation}
  C_{\mathrm{sample}} = \underbrace{O(d\cdot N_t)}_{\text{ODE}} + \underbrace{O(N_{\mathrm{obs}})}_{\text{interp.}} + \underbrace{O(n_p)}_{\text{prior}} + \underbrace{O(N_{\mathrm{obs}})}_{\text{likelihood}}.
\end{equation}
Total chain cost: $O\big((n_{\mathrm{samples}} + n_{\mathrm{burn}}) \cdot d \cdot N_t\big)$.

\subsection{Ising Model Monte Carlo}

Per sweep of an $L^d$ lattice:
\begin{align}
  \text{Single trial:}  &\quad O(1), \\
  \text{Full sweep:}    &\quad O(L^d), \\
  \text{Energy eval.:}  &\quad O(L^d) \;\text{(vectorised)}.
\end{align}

\subsection{GPU Batch Monte Carlo}

The GPU batch achieves near-linear weak scaling:
\begin{equation}
  T_{\mathrm{wall}} = O\!\left(N_t \cdot d \cdot \left\lceil \frac{B}{P}\right\rceil\right),
\end{equation}
where $B$ is the batch size and $P$ is the number of active GPU cores.  For $B \le P$ (typical for modern GPUs with $P \sim 10^3$--$10^4$ CUDA cores), wall-clock time is \emph{independent of batch size}.

%% =====================================================================
\section{Conclusion}\label{sec:conclusion}
%% =====================================================================

\pkg{MechanicsDSL} provides a cohesive MCMC architecture spanning three complementary subsystems.  The Bayesian inference engine couples the Metropolis--Hastings sampler directly to the symbolic--numeric simulation pipeline, enabling posterior inference over physical parameters with configurable priors, adaptive proposal scaling, and full chain diagnostics.  The Ising model implements a pedagogically rigorous single-site Metropolis simulator for 1D--3D lattice spin systems.  The CUDA batch code generator bridges the gap to high-performance computing by emitting embarrassingly parallel Monte Carlo kernels that exploit GPU hardware for massively concurrent parameter exploration.

Together, these subsystems enable a complete uncertainty quantification workflow: define a physical system in the DSL, derive equations of motion symbolically, fit parameters to data, sample the posterior, assess sensitivity, and propagate uncertainty---all within a unified framework.

%% =====================================================================
%% REFERENCES
%% =====================================================================
\begin{thebibliography}{9}

\bibitem{metropolis1953}
N.~Metropolis, A.\,W.~Rosenbluth, M.\,N.~Rosenbluth, A.\,H.~Teller, and E.~Teller,
``Equation of state calculations by fast computing machines,''
\textit{J.\ Chem.\ Phys.}, vol.~21, no.~6, pp.~1087--1092, 1953.

\bibitem{hastings1970}
W.\,K.~Hastings,
``Monte Carlo sampling methods using Markov chains and their applications,''
\textit{Biometrika}, vol.~57, no.~1, pp.~97--109, 1970.

\bibitem{roberts1997}
G.\,O.~Roberts, A.~Gelman, and W.\,R.~Gilks,
``Weak convergence and optimal scaling of random walk Metropolis algorithms,''
\textit{Ann.\ Appl.\ Probab.}, vol.~7, no.~1, pp.~110--120, 1997.

\bibitem{onsager1944}
L.~Onsager,
``Crystal statistics.\ I.\ A two-dimensional model with an order-disorder transition,''
\textit{Phys.\ Rev.}, vol.~65, no.~3--4, pp.~117--149, 1944.

\bibitem{saltelli2008}
A.~Saltelli \textit{et al.},
\textit{Global Sensitivity Analysis: The Primer}.
John Wiley \& Sons, 2008.

\end{thebibliography}

\end{document}
